% =====================================================================
%  AIMO Proof Pilot — Final grade template (problem captain)
%
%  HOW TO USE THIS FILE
%  --------------------
%  This is the FINAL grade we share with the teams. It is completed by
%  the problem captain once the independent grades have been reconciled.
%
%  1. COPY this file (do not edit it in place) and save it as
%        team_N_FINAL.tex
%     in the relevant problem's grades subfolder (e.g. grades/4_COLOUR/).
%  2. Fill in the final agreed mark and the reasoning behind it.
%  3. Compile with pdflatex.
%
%  Unlike the per-grader sheet, this records the single final decision:
%  the mark and the reasoning. It does not include the errors/issues log
%  or the qualitative proof-quality questions.
% =====================================================================
\documentclass[11pt]{article}
\usepackage[margin=2.2cm]{geometry}
\usepackage{amsmath}
\usepackage{amssymb}
\usepackage{booktabs}
\usepackage{array}
\usepackage{enumitem}
\usepackage{longtable}
\usepackage[hidelinks]{hyperref}
\usepackage{parskip}

\newcommand{\field}[1]{\textbf{#1:}\hspace{0.5em}}
% Placeholder for text the captain replaces. Renders as italic in [brackets]
% so it displays cleanly and is easy to spot. Delete it once filled in.
\newcommand{\fillin}[1]{{\itshape[#1]}}

\begin{document}

% ----------------------------------------------------------------------
%  Header — identify the submission
% ----------------------------------------------------------------------
{\centering\Large\bfseries AIMO Proof Pilot — Final Grade\par}
\vspace{1em}

\begin{tabular}{@{}p{0.5\textwidth}p{0.45\textwidth}@{}}
\field{Problem (\#\_ID)} \fillin{e.g. 4\_COLOUR}      & \field{Submission} \fillin{e.g. team\_3 / model\_2} \\
\end{tabular}

\vspace{0.4em}
\hrule
\vspace{0.8em}

% ----------------------------------------------------------------------
%  Reminder of the scoring scale
% ----------------------------------------------------------------------
\textbf{Scoring scale (0/1/6/7) — two binary decisions:}
\begin{enumerate}[label=\textbf{Step \arabic*.}, leftmargin=*, nosep]
    \item \textbf{Is the solution essentially complete?}
    \begin{itemize}[leftmargin=1.4em, nosep]
        \item If \textbf{no}: award \textbf{1} for significant partial progress, otherwise \textbf{0}.
        \item If \textbf{yes}: award \textbf{7}, or \textbf{6} if there is a minor error or omission.
    \end{itemize}
\end{enumerate}
Grade against the \textbf{markscheme} (\texttt{docs/markschemes.pdf}).

\vspace{0.8em}

% ----------------------------------------------------------------------
%  Final mark
% ----------------------------------------------------------------------
\section*{1. Final mark}

\textbf{Final mark} (choose one): \qquad \fbox{0} \qquad \fbox{1} \qquad \fbox{6} \qquad \fbox{7}

% ----------------------------------------------------------------------
%  Reasoning
% ----------------------------------------------------------------------
\section*{2. Reasoning}

Justify the mark with explicit reference to the markscheme points (e.g. \texttt{ME1}) where possible.

\bigskip
\fillin{Final reasoning here.}
\bigskip

\end{document}
